\subsubsection{Identifying Causal Effects}


\begin{Motivation}
    Consider an IL-BCN:
    \[ M = \lp G^+=(J,V,U,E^+),\quad \lp P_v\lp X_v|\doit(X_{\Pa^{G^+}(v)})   \rp\rp_{v\in V \cup U} \rp. \]
For simplicity assume that there are no input variables, i.e.\ $J=\emptyset$.
Then the joint distribution is ``$\doit$-free'' and given as:
\[P(X_V,X_U) = \prod_{v \in U\cup V}P_v\lp X_v|\doit(X_{\Pa^{G^+}(v)})\rp,   \]
with observational distribution as its marginal: $P(X_V)$.\\
We also have all the interventional distributions for $W \ins V$:
\[P(X_{V \sm W},X_U|\doit(X_W)) = \prod_{v \in U \cup V \sm W}P_v\lp X_v|\doit(X_{\Pa^{G^+}(v)})\rp,\]
with marginals: $P(X_{V\sm W}|\doit(X_W))$.\\
If we wanted to learn the distribution $P(X_V)$ we could do an observational study and apply the usual statistical or machine learning techniques.
If, in contrast, we wanted to learn interventional distributions: $P(X_{V\sm W}|\doit(X_W))$ from data, e.g.\ 
if vaccination makes people immune to a disease,
we typically would need to perform an interventional study where we intervene on the variables $X_W$ on different values.
This usually requires expensive, time-consuming randomized control trials with an own comparison group for each possible 
values of $X_W$.\\
If we assume that we know the causal graph $G^+$ or $G$ we could try to leverage the rules of do-calculus in a clever way and might be able to go from expressions involving $\doit(W)$ to expressions only involving $\doit(D)$ 
for a (much) smaller subset $D \ins W$, ideally $D=\emptyset$.\\
Practically this would mean that we would need a much smaller randomized control trial and save time and resources.\\
For example, if we have the graph only involving the edge: $v_1 \tuh v_2$ we have that:
\[ P(X_2|\doit(X_1)) =  P(X_2|X_1), \]
which can be estimated using observational data only, e.g.\ via supervised learning.\\
So the question is now: 
Assuming that the causal graph is known, under which circumstances is a causal effect 
$P(X_A|\doit(X_B))$ already determined by the observational distribution $P(X_V)$?
When can causal effects be identified via distributions that have less interventions in them?
\end{Motivation}

\begin{Not}
Consider an IL-BCN:
    \[ M = \lp G^+=(J,V,U,E^+),\quad \lp P_v\lp X_v|\doit(X_{\Pa^{G^+}(v)})   \rp\rp_{v\in V \cup U} \rp. \]
We are interested in estimating the conditional causal effect:
\[ P(X_A|X_C,\doit(X_B,X_D)), \]
but we only have data from:
\[ P(X_V|X_C,\doit(X_D)). \]
The following index sets will have the following roles:
\begin{enumerate}
    \item $A$: the outcome variables of interest.
    \item $B$: the treatment or intervention variables.
    \item $C$: general conditional (context) variables under which the data was collected.
    \item $D$: general interventional (context) variables that were set by the experimenter, $J \ins D$.
    \item $F_0$: core adjustment variables, i.e.\ features that were measured.
    \item $F_1$: additional measured adjustment variables.
    \item $F = F_0 \cup F_1$.
    \item $H$: additional unobserved variables.
\end{enumerate}


\end{Not}


\begin{Thm}[General adjustment formula]
    \label{thm-gen-adjustment}
Consider an IL-BCN:
    \[ M = \lp G^+=(J,V,U,E^+),\quad \lp P_v\lp X_v|\doit(X_{\Pa^{G^+}(v)})   \rp\rp_{v\in V \cup U} \rp. \]
    Assume that all the following conditions hold in the graphs $G^+_{\doit(I_B,D)}$: % $G^+_{\doit(D)}$:
    \begin{align}
        (F_0 \cup H) & \Perp^d_{G^+_{\doit(I_B,D)}} I_B \given (C \cup D), \label{eq:adj-1} \\
%        H & \Perp^d_{G^+_{\doit(I_B,D)}} I_B \given (C \cup D), \label{eq:adj-2} \\
        A & \Perp^d_{G^+_{\doit(I_B,D)}} (F_1 \cup I_B) \given (B \cup F_0 \cup H \cup C \cup D),\label{eq:adj-3} \\
%        A & \Perp^d_{G^+_{\doit(I_B,D)}} F_1 \given ( B \cup F_0 \cup H \cup C \cup D),\label{eq:adj-4} \\
                      H &\Perp^d_{G^+_{\doit(I_B,D)}} B \given (F \cup C \cup  I_B \cup D). \label{eq:adj-5} 
    \end{align}
Then we have the adjustment formula:
\[ P(X_A|X_C,\doit(X_B,X_D)) = 
    P(X_A|X_B,X_C,X_F,\doit(X_D)) \circ P(X_F|X_C,\doit(X_D)) \qquad \text{a.s.}
\]
\begin{proof}
%    By assumption \ref{eq:adj-3}:
%    \[ A  \Perp^d_{G^+_{\doit(I_B,D)}} (F_1 \cup I_B) \given (B \cup F_0 \cup H \cup C \cup D),\]
%there is a Markov kernel that simultaneously is a version of:
%\[P(X_A|X_{F_0},X_H,X_C,\doit(X_B,X_D)),\qquad 
%P(X_A|X_{F_0},X_H,X_C,X_B,\doit(X_D)),\]
%and:
%\[P(X_A|X_{F_0},X_{F_1},X_H,X_C,X_B,\doit(X_D)).\]
%We denote it with:
%\[ P(X_A|X_{F_0},\cancel{X_{F_1}},X_H,X_C,\cancel{\doit}(X_B),\doit(X_D))\]
    \begin{align}
 & P(X_A|X_C,\doit(X_B,X_D)) \\
 &= P(X_A|X_{F_0},X_H,X_C,\doit(X_B,X_D)) \circ P(X_{F_0},X_H|X_C,\doit(X_B,X_D))\\
 &\stackrel{\ref{eq:adj-3}}{=}
 P(X_A|X_{F_0},X_H,X_C,X_B,\doit(X_D))\circ P(X_{F_0},X_H|X_C,\doit(X_B,X_D)) \\
 &\stackrel{\ref{eq:adj-1}}{=}
 P(X_A|X_{F_0},X_H,X_C,X_B,\doit(X_D))\circ P(X_{F_0},X_H|X_C,\doit(X_D)) \\
 &\stackrel{}{=}
 P(X_A|X_{F_0},X_H,X_C,X_B,\doit(X_D))\circ P(X_{F_0},X_{F_1},X_H|X_C,\doit(X_D)) \\
 & \stackrel{\ref{eq:adj-3}}{=}
 P(X_A|X_{F_0},X_{F_1},X_H,X_C,X_B,\doit(X_D))\circ P(X_{F_0},X_{F_1},X_H|X_C,\doit(X_D)) \\
  & \stackrel{F=F_0\dcup F_1}{=}
 P(X_A|X_F,X_H,X_C,X_B,\doit(X_D))\circ P(X_F,X_H|X_C,\doit(X_D)) \\
 & \stackrel{}{=}
 P(X_A|X_F,X_H,X_C,X_B,\doit(X_D))\circ \lp P(X_H|X_F,X_C,\doit(X_D)) \otimes P(X_F|X_C,\doit(X_D))  \rp \\
  & \stackrel{\ref{eq:adj-5}}{=}
 P(X_A|X_F,X_H,X_C,X_B,\doit(X_D))\circ \lp P(X_H|X_F,X_C,X_B,\doit(X_D)) \otimes P(X_F|X_C,\doit(X_D))  \rp \\
   & \stackrel{}{=}
 P(X_A|X_F,X_C,X_B,\doit(X_D))\circ  P(X_F|X_C,\doit(X_D)) \\
    & \stackrel{}{=}
 P(X_A|X_B,X_F,X_C,\doit(X_D))\circ  P(X_F|X_C,\doit(X_D)).
    \end{align}
\end{proof}
\end{Thm}


\begin{Cor}[Conditional backdoor covariate adjustment formula]
    \label{thm-backdoor-cond-int}
    Consider an IL-BCN:
    \[ M = \lp G^+=(J,V,U,E^+),\quad \lp P_v\lp X_v|\doit(X_{\Pa^{G^+}(v)})   \rp\rp_{v\in V \cup U} \rp. \]
  Assume that the \emph{conditional backdoor criterion} in the graphs $G^+_{\doit(I_B,D)}$ holds:
\begin{enumerate}
    \item $\displaystyle F \Perp^d_{G_{\doit(I_B,D)}} I_B  \given (C \cup D)$, and:
    \item $\displaystyle A \Perp^d_{G_{\doit(I_B,D)}} I_B  \given (B \cup F \cup C \cup D)$.
\end{enumerate}
Then we have the adjustment formula:
\[ P(X_A|X_C,\doit(X_B,X_D)) = 
    P(X_A|X_B,X_F,X_C,\doit(X_D)) \circ P(X_F|X_C,\doit(X_D)) \qquad \text{a.s.}
\]
\begin{proof}
    It follows directly from theorem \ref{thm-gen-adjustment} by using $F_1=H=\emptyset$.\\
    We give another proof in more details here anyways.\\
    By the assumption:
 \[ A \Perp^d_{G_{\doit(I_B,D)}} I_B  \given (B \cup F \cup C \cup D),\]
 there exists a Markov kernel that is simultaneously a version of:
 \[ P(X_A|X_F,X_C, \doit(X_B,X_D)) \qquad\text{and}\qquad P(X_A|X_F,X_C, X_B,\doit(X_D)).\]
 We denote it by:
 \[ P(X_A|X_F,X_C, \cancel{\doit}(X_B),\doit(X_D)). \]
 By assumption:
 \[F \Perp^d_{G_{\doit(I_B,D)}} I_B  \given (C \cup D),\]
 there exists a Markov kernel that is simultaneously a version of:
 \[ P(X_F|X_C, \doit(X_B,X_D)) \qquad\text{and}\qquad P(X_F|X_C,\doit(X_D)).\]
 We denote it by:
 \[ P(X_F|X_C,\cancel{\doit(X_B)},\doit(X_D)). \]
 Then:
 \[P(X_A|X_F,X_C, \cancel{\doit}(X_B),\doit(X_D)) \circ P(X_F|X_C,\cancel{\doit(X_B)},\doit(X_D)) \]
 is a version of:
 \[ P(X_A|X_C,\doit(X_B,X_D)).\]
\end{proof}
\end{Cor}


\begin{Cor}[Backdoor covariate adjustment]\label{cor:backdoor_adjustment}
    Let the situation be like in theorem \ref{thm-backdoor-cond-int} with $C=D=J= \emptyset$.
    Assume that the \emph{backdoor criterion} holds:
\begin{enumerate}
    \item $\displaystyle F \Perp^d_{G_{\doit(I_B)}} I_B$, and:
    \item $\displaystyle A \Perp^d_{G_{\doit(I_B)}} I_B  \given (B \cup F)$.
\end{enumerate}
Then we have the adjustment formula:
\[ P(X_A|\doit(X_B)) = 
    P(X_A|X_B,X_F) \circ P(X_F) \qquad \text{a.s.}
\]
\end{Cor}

\begin{Rem}[Tian's ID algorithm]
    There exists an algorithm that takes as input the observational CADMG $G$ of an IL-CBN that have densities/mass functions,
    and sets: $A,B,C \ins J \cup V$.
    It outputs either an algebraic formula for $P(X_A|X_B,\doit(X_C))$ in terms of conditional marginals 
    of the observational Markov kernel $P(X_V|\doit(X_J))$ if the above causal effect is identifiable, or: it outputs FAIL 
    if the above causal effect is not identifiable.
    The algorithm only uses the do-calculus rules.
\end{Rem}
\Joris{Patrick suggested that the scope of this claim may need to be restricted somehow.}
